\section{两种正交变换}

我们知道，正交矩阵$\vb{Q}\in\R^{m\times m}$是指满足以下条件的矩阵
\begin{Equation}[正交矩阵定义]
    \vb{Q}^T\vb{Q}=\vb{Q}\vb{Q}^T=\vb{I}
\end{Equation}

在LU分解中，消元依靠的是高斯变换。在QR分解中，消元依靠的则是正交变换！可以这样理解：QR分解同样要产生上三角矩阵$\vb{R}$，该过程就是由一系列正交变换构成的$\vb{Q}$实现的。

在$\R^{2\times 2}$下，有两种重要的正交变换：反射和旋转。对于$\theta$，我们记
\begin{Equation}[2x2正交cs]
    c=\cos\theta\qquad s=\sin\theta
\end{Equation}

旋转（Rotation）矩阵$\vb{Q}$可以表示为以下形式，代表旋转$\theta$的角度
\begin{Equation}[2x2正交旋转]
    \vb{Q}=\mqty[
        c & s \\
        -s & c \\
    ]
\end{Equation}

反射（Reflection）矩阵$\vb{Q}$可以表示为以下形式，代表沿着$\theta/2$反射
\begin{Equation}[2x2正交反射]
    \vb{Q}=\mqty[
        c & s \\
        s & -c \\
    ]
\end{Equation}

这一节要研究的问题是，如何将$\R^{2\times 2}$上的旋转和反射以某种方式应用到$\R^{m\times m}$的矩阵上。

\subsection{Householder反射变换}

在$\R^m$下，通常反射会选取一个法向量$\vb{v}\in\R^m$并讨论$\vb{x}\in\R^m$对于垂直$\vb{v}$的超平面的镜像。

\begin{BoxFormula}[Householder反射变换]
    Householder反射变换矩阵$\vb{H}\in\R^{m\times m}$可以表示为以下形式
    \begin{Equation}[Householder反射变换]
        \vb{H}(\vb{v})=\vb{I}-\beta\vb{v}\vb{v}^T\qquad\beta=\frac{2}{\vb{v}^T\vb{v}}
    \end{Equation}
    其中$\vb{v}\in\R^m$为反射超平面的法向量。
\end{BoxFormula}

现在的问题是，对于任意的$\vb{x}$，如何找到对应的$\vb{v}$构造出变换矩阵$\vb{H}$，使$\vb{H}\vb{x}=\vb{e}_1$成立，换言之，这就是将$\vb{x}$变换为只有第一个元素为$1$其余元素都为$0$的单位向量$\vb{e}_1$，即所谓消元！
\begin{Equation}[反射变换1]
    \mqty[
        \vb{x}(1)\\
        \vb{x}(2)\\
        \vdots\\
        \vb{x}(m)\\
    ]\to\mqty[
        \times\\
        0\\
        \vdots\\
        0
    ]
\end{Equation}

根据\cref{fml:Householder反射变换}可得，请注意由于$\vb{v}^T\vb{x}$是数，所以可以提到前面
\begin{Equation}[反射变换2]
    \vb{H}\vb{x}=\qty(\vb{I}-\frac{2\vb{v}\vb{v}^T}{\vb{v}^T\vb{v}})\vb{x}=\vb{x}-\frac{2\vb{v}^T\vb{x}}{\vb{v}^T\vb{v}}\vb{v}
\end{Equation}

由于$\vb{v},\vb{x},\vb{e}_1$的反射关系，三者共平面，故可以将$\vb{v}$写为$\vb{x},\vb{e}_1$的线性组合
\begin{Equation}[反射变换3]
    \vb{v}=\vb{x}+\alpha\vb{e}_1
\end{Equation}

注意到\cref{eq:反射变换2}的分子和分母有$\vb{v}^T\vb{x}$和$\vb{v}^T\vb{v}$，现在要用\cref{eq:反射变换3}表示这两者。

项$\vb{v}^T\vb{x}$可以表示为
\begin{Equation}[反射变换4]
    \vb{v}^T\vb{x}=\vb{x}^T\vb{x}+\alpha\vb{x}(1)
\end{Equation}

项$\vb{v}^T\vb{v}$可以表示为
\begin{Equation}[反射变换5]
    \vb{v}^T\vb{v}=\vb{x}^T\vb{x}+2\alpha\vb{x}(1)+\alpha^2
\end{Equation}

现在将\cref{eq:反射变换3,eq:反射变换4,eq:反射变换5}代入\cref{eq:反射变换2}
\begin{Equation}[反射变换6]
    \vb{H}\vb{x}=1-2\frac{\vb{x}^T\vb{x}+\alpha\vb{x}(1)}{\vb{x}^T\vb{x}+2\alpha\vb{x}(1)+\alpha^2}\qty(\vb{x}+\alpha\vb{e}_1)
\end{Equation}

稍做整理
\begin{Equation}[反射变换7]
    \vb{H}\vb{x}=\qty(1-2\frac{\vb{x}^T\vb{x}+\alpha\vb{x}(1)}{\vb{x}^T\vb{x}+2\alpha\vb{x}(1)+\alpha^2})\vb{x}-2\alpha\frac{\vb{v}^T\vb{v}}{\vb{v}^T\vb{v}}\vb{e}_1
\end{Equation}

化简
\begin{Equation}[反射变换8]
    \vb{H}\vb{x}=\qty(\frac{\alpha^2-\vb{x}^T\vb{x}}{\vb{x}^T\vb{x}+2\alpha\vb{x}(1)+\alpha^2})\vb{x}-2\alpha\frac{\vb{v}^T\vb{v}}{\vb{v}^T\vb{v}}\vb{e}_1
\end{Equation}

请注意$\vb{v}^T\vb{v}$实际上就是$\norm{\vb{v}}_2^2$
\begin{Equation}[反射变换9]
    \vb{H}\vb{x}=\qty(\frac{\alpha^2-\norm{\vb{v}}_2^2}{\vb{x}^T\vb{x}+2\alpha\vb{x}(1)+\alpha^2})\vb{x}-2\alpha\frac{\vb{v}^T\vb{v}}{\vb{v}^T\vb{v}}\vb{e}_1
\end{Equation}

我们注意到，如果此时令$\alpha=\pm\norm{\vb{x}}_2$，就有$\vb{v}=\vb{x}\pm\norm{\vb{x}}_2\vb{e}$以及
\begin{Equation}[反射变换10]
    \vb{H}\vb{x}=-2\alpha\frac{\vb{v}^T\vb{v}}{\vb{v}^T\vb{v}}\vb{e}_1
\end{Equation}

现在已经知道$\vb{H}\vb{x}$平行于$\vb{e}_1$了，但到底$\vb{H}\vb{x}$是多少倍的$\vb{e}_1$呢？代入\cref{eq:反射变换4,eq:反射变换5}
\begin{Equation}[反射变换11]
    \vb{H}\vb{x}=-2\alpha\frac{\vb{x}^T\vb{x}+\alpha\vb{x}(1)}{\vb{x}^T\vb{x}+2\alpha\vb{x}(1)+\alpha^2}\vb{e}_1
\end{Equation}

代入$\alpha=\pm\norm{\vb{x}}_2$和$\vb{x}^T\vb{x}=\norm{\vb{x}}_2^2$
\begin{Equation}[反射变换12]
    \vb{H}\vb{x}=\mp\norm{\vb{x}}_2\vb{e}_1
\end{Equation}

我们将这个结论整理如下
\begin{BoxFormula}[Householder矩阵的构造]
    对于$\vb{x}\in\R^m$，若选取以下Householder向量$\vb{v}\in\R^m$
    \begin{Equation}[Householder矩阵的构造1]
        \vb{v}=\vb{x}\pm\norm{\vb{x}}_2\vb{e}_1
    \end{Equation}
    则$\vb{v}$构造的Householder矩阵满足
    \begin{Equation}[Householder矩阵的构造2]
        \vb{H}\vb{x}=\mp\norm{\vb{x}}_2\vb{e}_1
    \end{Equation}
\end{BoxFormula}

如\cref{fig:Householder反射变换}所示，请注意$\vb{x}$是关于垂直于$\vb{v}$的超平面（而不是关于$\vb{v}$）进行镜像的。

\begin{Figure}[Householder反射变换]
    \includegraphics[scale=0.9]{HouseholderReflect.fig.pdf}
\end{Figure}

通常来说我们会选取$\vb{v}=\vb{x}-\norm{\vb{x}}_2\vb{e}_1$以便得到$\vb{H}\vb{x}=\norm{\vb{x}}_2\vb{e}_1$的正的变换结果。

\begin{BoxAlgorithm}[Householder参数]
    对于$\vb{x}\in\R^m$，该函数计算满足$\vb{v}(1)=1$的$\vb{v}\in\R^m$和$\beta\in\R$的Householder参数。
    
    使$\vb{H}=\vb{I}-\beta\vb{v}\vb{v}^T$是Householder矩阵，且满足$\vb{H}\vb{x}=\norm{\vb{x}}_2\vb{e}_1$。
    \begin{Algoplain}
        \Function{$[v,\beta]=\mathrm{house}(x)$}{
            $\vb{v}(1)=1,\ \vb{v}(2:m)=\vb{x}(2:m),\ \sigma=\vb{x}(2:m)\vb{x}(2:m)^T$\;
            \uIf{$\sigma=0$ and $\vb{x}(1)\geq 0$}{$\beta=0$\;}
            \uElseIf{$\sigma=0$ and $\vb{x}(1)< 0$}{$\beta=-2$\;}
            \Else{
                $\mu=\sqrt{\vb{x}(1)^2+\sigma}$\;
                \uIf{$\vb{x}(1)\leq 0$}{$\vb{v}(1)=\vb{x}(1)-\mu$}   
                \Else{$\vb{v}(1)=-\sigma/(\vb{x}(1)+\mu)$}
                $\beta=2\vb{v}(1)^2/(\sigma+\vb{v}(1)^2)$\;
                $\vb{v}=\vb{v}/\vb{v}(1)$\;
            }
        }
    \end{Algoplain}
\end{BoxAlgorithm}

然而，在算法层面，直接计算$\vb{v}=\vb{x}-\norm{\vb{x}}_2\vb{e}_1$有一些额外的麻烦，注意到
\begin{Equation}[Householder参数1]
    \vb{v}(2:m)=\vb{x}(2:m)
\end{Equation}

因此其实真正要关心的只有$\vb{v}(1)$，我们会分为$\vb{x}(1)\leq 0$和$\vb{x}(1)>0$两种情况来讨论。

对于$\vb{x}(1)\leq 0$，直接进行计算
\begin{Equation}[Householder参数2]
    \vb{v}(1)=\vb{x}(1)-\norm{\vb{x}}_2
\end{Equation}

对于$\vb{x}(1)>0$，考虑进行如下变换以避免数值精度问题
\begin{Equation}[Householder参数3]
    \vb{v}(1)=\frac{\vb{x}(1)-\norm{\vb{x}}_2^2}{\vb{x}(1)+\norm{\vb{x}}_2}=-\frac{\vb{x}(2)^2+\cdots+\vb{x}(m)^2}{\vb{x}(1)+\norm{\vb{x}}_2}
\end{Equation}

引入代换变量$\sigma$和$\mu$，并注意到
\begin{Equation}[Householder参数4]
    \sigma=\vb{x}(2)^2+\cdots+\vb{x}(m)^2\qquad
    \mu=\sqrt{\vb{x}(1)+\sigma}=\norm{\vb{x}}_2
\end{Equation}

这样一来，\cref{eq:Householder参数2,eq:Householder参数3}就可以表示为
\begin{Equation}[Householder参数5]
    \vb{v}(1)=\begin{cases}
        \vb{x}(1)-\mu,&\vb{x}(1)\leq 0\\
        -\sigma/(\vb{x}(1)+\mu),&\vb{x}(1)>0\\
    \end{cases}
\end{Equation}

除此之外，在\cref{alg:Householder参数}中，还有一些值得注意的细节。算法期望缩放$\vb{v}$使$\vb{v}(1)=1$，目的是存储$\vb{v}$时可以少用一位，这固然可行，但这种缩放一定会影响$\vb{H}=\vb{I}-\beta\vb{v}\vb{v}^T$的结果。由于$\beta$也由算法给出，因此在计算$\beta$时额外引入了$\vb{v}(1)^2$以补偿稍后$\vb{v}=\vb{v}/\vb{v}(1)$在$\beta\vb{v}\vb{v}^T$上的损失。

\subsection{Givens旋转变换}

在$\R^m$下，通常旋转仍是对其中特定的两个元素进行，需要将\cref{eq:2x2正交旋转}放到$\R^{m\times m}$中完成。

\begin{BoxFormula}[Givens旋转变换]
    Givens旋转变换矩阵$\vb{G}\in\R^{m\times m}$可以表示为以下形式
    \begin{Equation}[Givens旋转变换]
        \vb{G}(i,j,\theta)=
        \begin{bNiceMatrix}[first-row, last-col]
                    &i      &       &j      &       &       \\
            \vb{I}  &       &       &       &       &       \\
                    &+c     &       &+s     &       &i      \\
                    &       &\vb{I} &       &       &       \\
                    &-s     &       &+c     &       &j      \\
                    &       &       &       &\vb{I} &       \\
        \end{bNiceMatrix}
    \end{Equation}
    其中$\theta$是旋转角（$c=\cos\theta,\ s=\sin\theta$），而$i,j$是将要旋转的元素下标。
\end{BoxFormula}

现在我们想要达成的是，对于任意的$\vb{x}$，通过对$\vb{x}(i)$和$\vb{x}(j)$的旋转使后者归零。
\begin{Equation}[旋转变换1]
    \mqty[
        \vb{x}(1)\\
        \vdots\\
        \vb{x}(i)\\
        \vdots\\
        \vb{x}(j)\\
        \vdots\\
        \vb{x}(m)
    ]\to
    \mqty[
        \vb{x}(1)\\
        \vdots\\
        \times\\
        \vdots\\
        0\\
        \vdots\\
        \vb{x}(m)
    ]
\end{Equation}

请注意Givens矩阵$\vb{G}$表旋转示$\theta$时需要$\vb{G}^T$左乘，不妨记
\begin{Equation}[旋转变换2]
    \vb{G}(i,k,\theta)^T\vb{x}=\vb{y}
\end{Equation}

注意到$\vb{y}(k)$满足
\begin{Equation}[旋转变换3]
    \vb{y}(k)=\begin{cases}
        c\vb{x}(i)-s\vb{x}(j),&k=i\\
        s\vb{x}(i)+c\vb{x}(j),&k=j\\
        \vb{x}(k),&k\neq i,j\\
    \end{cases}
\end{Equation}

通过\cref{eq:旋转变换3}，可以确定，如果要令$\vb{y}(j)=0$，应当将$c,s$取为
\begin{Equation}[旋转变换4]
    c=\frac{\vb{x}(i)}{\sqrt{\vb{x}(i)^2+\vb{x}(j)^2}}\qquad
    s=\frac{-\vb{x}(j)}{\vb{x}(i)^2+\vb{x}(j)^2}
\end{Equation}

我们可以将结论整理如下
\begin{BoxFormula}[Givens矩阵的构造]
    对于$\vb{x}\in\R^m$，若选取以下Givens参数$c,s$
    \begin{Equation}[Givens矩阵的构造1]
        c=\frac{\vb{x}(i)}{\sqrt{\vb{x}(i)^2+\vb{x}(j)^2}}\qquad
        s=\frac{-\vb{x}(j)}{\sqrt{\vb{x}(i)^2+\vb{x}(j)^2}}
    \end{Equation}
    并记$c,s$构造的Givens矩阵作用在$\vb{x}$上的变换结果为$\vb{y}$
    \begin{Equation}[Givens矩阵的构造2]
        \vb{G}^T\vb{x}=\vb{y}
    \end{Equation}
    则其满足
    \begin{Equation}[Givens矩阵的构造3]
        \vb{y}(j)=0
    \end{Equation}
\end{BoxFormula}

在算法层面，只关注需要旋转的元素，且有一些技巧可以减少运算并提高数值精度。

\begin{BoxAlgorithm}[Givens参数]
    对于$\vb{x}\in\R^2$，该函数计算$c\in\R$和$s\in\R$的Givens参数，并满足
    \begin{Equation}[Givens参数算法]
        \mqty[
            c & s \\ 
            -s & c \\
        ]^T
        \mqty[
            a\\
            b\\
        ]=
        \mqty[
            r\\
            0\\
        ]
    \end{Equation}
    \begin{Algoplain}
        \Function{$[c,s]=\mathrm{givens}(x)$}{
            \uIf{$b=0$}{$c=1,\ s=0$}
            \Else{
                \uIf{$\abs{b}>\abs{a}$}{$\tau=-a/b,\ s=1/\sqrt{1+\tau^2},\ c=s\tau$\;}
                \Else{$\tau=-b/a,\ c=1/\sqrt{1+\tau^2},\ s=c\tau$\;}
            }
        }
    \end{Algoplain}
\end{BoxAlgorithm}
